\scanpage{53}
\mspage{\handwritten{- 12 -}}
% archive sheet number: 52
En effet, en désignant par $\handwritten{\unsure{c|C}}_{k;n_1,\ldotp\ldotp,n_p}$ les coefficients de $\handwritten{f}_k$, les coefficients $\handwritten{b}_{n_1\ldotp\ldotp n_p}$ de $\handwritten{g}$ (supposée exister) doivent satisfaire aux conditions $n_k\handwritten{b}_{n_1\ldotp\ldotp n_p}=\handwritten{c}_{k;n_1,\ldotp\ldotp,n_k-1,\ldotp\ldotp,n_p}$ pour $(n_1,\ldotp\ldotp,n_p)\handwritten{\neq}\handwritten{\neq}(0,\ldotp\ldotp,0)$ et $\handwritten{b}_{00\ldotp\ldotp0}=\handwritten{c}$ ; l'hypothèse implique que ces conditions déterminent les $\handwritten{b}_{n_1\ldotp\ldotp n_p}$ de façon unique. En outre, on a, dans $K$, $\handwritten{\bigl|}\frac{1}{n}\handwritten{\bigr|}\handwritten{\leq}n^{\handwritten{\beta}}$ pour une constante fixe $\handwritten{\beta>}0$ ne dépendant que de la valeur absolue de $K$, et on vérifie aisément que dans $\handwritten{\mathbb{R}}$, la série de terme général
\[
((n_1+1)(n_2+1)\ldotp\ldotp(n_p+1))^{\handwritten{\beta}}\handwritten{\|}\handwritten{a}_{n_1n_2\ldotp\ldotp n_p}\handwritten{\|}r_1^{n_1}r_2^{n_2}\ldots r_p^{n_p}
\]
converge en même temps que la série de terme général
\[
\handwritten{\|}\handwritten{a}_{n_1\ldotp\ldotp n_p}\handwritten{\|}r_1^{n_1}r_2^{n_2}\ldots r_p^{n_p}.
\]

PROPOSITION 8.- \uline{Soit} $\handwritten{f}(z_1,\ldotp\ldotp,z_p)$ \uline{une fonction holomorphe dans un ensemble ouvert} $A$. \uline{Pour tout point} $\handwritten{a}=(a_1,\ldotp\ldotp,a_p)$, \uline{de} $A$ \uline{et tout nombre} $\handwritten{\varepsilon>}0$, \uline{il existe un polycylindre ouvert} $P_o$ \uline{de centre} $\handwritten{a}$ \uline{tel que}, \uline{pour tout couple de points} $(u_1,\ldotp\ldotp,u_p),(v_1,\ldotp\ldotp,v_p)$ \uline{de ce polycylindre}, \uline{on ait}
\[
\handwritten{\Bigl\|}\handwritten{f}(u_1,\ldotp\ldotp,u_p)-\handwritten{f}(v_1,\ldotp\ldotp,v_p)-\handwritten{\sum}_{\handwritten{k=1}}^{p}D_k\handwritten{f}(a_1,\ldotp\ldotp,a_p)(u_k-v_k)\handwritten{\Bigr\|}\handwritten{\leq}\handwritten{\varepsilon}.\sup_{1\handwritten{\leq}k\handwritten{\leq}p}\handwritten{|}u_k-v_k\handwritten{|}
\tag*{(6)}
\]

On peut se ramener au cas où $\handwritten{a}$ est l'origine, $\handwritten{f}(z_1,\ldotp\ldotp,z_p)$ étant la somme d'une série entière $(\handwritten{c}_{n_1\ldotp\ldotp n_p}z_1^{n_1}\ldotp\ldotp z_p^{n_p})$ absolument sommable dans un polycylindre ouvert $Q$. En vertu du th.2, le polynome $u_k+w_k$ est substituable à $z_k$ dans cette série entière (pour $1\handwritten{\leq}k\handwritten{\leq}p$), dès que $(u_k)$ et $(w_k)$ appartiennent à un polycylindre $P$ assez petit, et par substitution, on obtient une série entière absolument sommable dans $P\handwritten{\times}P$ par rapport à ces $2p$ variables. Par le théorème d'associativité, on peut donc écrire
